\section{视觉语言接口：从 Q-Former 到线性投影}

把 image 接入 LLM 的第一个问题不是“模型看懂了吗”，而是视觉 encoder 的连续特征怎样进入 language model 的 token space。BLIP-2 与 LLaVA 给出两种典型桥接方式：一个使用可学习 queries 和 Q-Former 压缩/对齐，一个直接线性投影视觉 token。

\subsection{BLIP-2：用 Q-Former 查询冻结视觉空间}

BLIP-2 把 frozen image encoder 与 frozen LLM 之间放入 Q-Former。令视觉特征为 $V\in\mathbb{R}^{m\times d_v}$，可学习 queries 为 $Q_0\in\mathbb{R}^{q\times d}$，cross attention 为
\[
Q'=\operatorname{softmax}\!\left(\frac{(Q_0W_Q)(VW_K)^\top}{\sqrt d}\right)VW_V.
\]
$q\ll m$ 时，Q-Former 同时做信息瓶颈与空间对齐，再把少量 query outputs 送入 LLM。

\lecturefigure{slide-13-blip2.jpg}{BLIP-2：Q-Former 桥接 frozen image encoder 与 frozen LLM}{Stanford Online 官方录像 00:36:18}

\begin{knowledgebox}{Q-Former 解决两个不同问题}
一是\textbf{compression}：从大量 image patches 提炼固定数量 queries；二是\textbf{alignment}：把 CLIP-like feature space 转到 LLM 可消费的 embedding space。它不是让 LLM 自己重新学习全部视觉 backbone。
\end{knowledgebox}

\teachervoice{00:35:00--00:36:55，讲者把 BLIP-2 称为有效连接 CLIP 与 LLM 的先驱，重点在“两个冻结空间之间需要可训练桥”。}

\subsection{LLaVA：简单 projection 为什么成为强 baseline}

LLaVA 使用 vision encoder $f_v$ 和 projection $W_p$：
\[
H_v=f_v(I),\qquad Z_v=H_vW_p,
\]
再把 $Z_v$ 与 text embeddings 拼接输入 LLM。它减少了额外模块，容易 end-to-end instruction tuning，因此很快成为 VLM 常见 baseline。

\lecturefigure{slide-14-llava.jpg}{LLaVA：以 projection 对齐视觉特征与语言 embedding}{Stanford Online 官方录像 00:36:55}

\begin{warningbox}{Linear projection 简单，不等于信息无损}
projection 只保证维度匹配。视觉 encoder 的分辨率、patch size、pretraining objective、projector capacity 和 instruction data 决定哪些细节能进入 LLM；小文字、OCR、grounding 与 GUI icons 往往首先暴露瓶颈。
\end{warningbox}

\subsection{接口 taxonomy：压缩、对齐、保留与高分辨率}

在进入 Cog 系列前，可以用四个问题比较所有 VLM bridge：多少视觉 token 被保留？视觉与语言是否共享参数？是否会破坏原语言能力？高分辨率输入的 cost 在哪里支付？

\begin{knowledgebox}{VLM bridge 对照表}
\begin{tabularx}{\textwidth}{L{0.16\textwidth} L{0.22\textwidth} L{0.23\textwidth} X}
\toprule
方法 & 视觉接口 & 参数路径 & 主要权衡 \\
\midrule
BLIP-2 & Q-Former queries & bridge 可训练，backbones 可冻结 & 压缩强，但 query bottleneck 明显 \\
LLaVA & linear/MLP projector & 视觉 token 进入 LLM path & 简单强基线，高分辨率成本高 \\
CogVLM & vision experts & image/text token 走部分不同参数 & 尝试保留语言行为，参数更多 \\
CogAgent & high-res cross attention & 轻量高分辨率旁路 & 面向 GUI/OCR，系统更复杂 \\
Vary & 多视觉词表/feature ensemble & 多 encoder 特征合并 & OCR 更强，接口与训练成本上升 \\
GLM-4V slide & stride convolution & 压缩后 mixed training & 结构简单，结果依赖数据与 protocol \\
\bottomrule
\end{tabularx}
\end{knowledgebox}

\subsection{本章小结}

多模态接口的本质是 information bottleneck 与 parameter sharing。Q-Former 显式查询和压缩，LLaVA 直接投影；后续模型则围绕“保留语言能力”和“处理高分辨率细节”继续分化。

